% Werke X/1 (1917), native237–238, printed230–231.
\begingroup
\clearpage\SourceLeaf{237}\ArithmeticHead{230}{L}{ANALYSIS.\quad NACHLASS.}
\noindent Dann ist, den letzten Wert von \(T=\theta\) gesetzt,
{\small
\[\begin{aligned}
W&=\frac\theta\alpha(a'a'-32c''c''-64c'''c'''-128c^{\mathrm{IV}}c^{\mathrm{IV}}\ldots)\\
&\quad+8c'\cos T\sin T'+16c''\cos T'\sin T''+32c'''\cos T''\sin T'''+\ldots\\[.8em]
&=\frac\theta\alpha(a'a'-32c''c''-64c'''c'''-128c^{\mathrm{IV}}c^{\mathrm{IV}}\ldots)\\
&\quad+2a'p'\cos T+4a''p''\cos T'+8a'''p'''\cos T''\ldots
\end{aligned}\]}
\par\indent Salv.\ factore const.
{\small\setlength{\arraycolsep}{6pt}
\[\begin{array}{r@{\;=\;}r@{}c@{}r|l}
a&302{,}78000&\quad&2{,}4811272&c\;\ldots\;0{,}7886597\\[.4em]
b&301{,}78000&\quad&2{,}4796904&c'\;\ldots\;9{,}0969100-10\\[.2em]\cline{1-2}\cline{4-4}\noalign{\vskip.2em}
a'&302{,}28000&\quad&2{,}4804094&c''\;\ldots\;5{,}7134110-10\\[.4em]
b'&302{,}27944&\quad&2{,}4804086&c'''\;\ldots\;8{,}9464130-20\\[.2em]\cline{1-2}\cline{4-4}\noalign{\vskip.2em}
\alpha&302{,}27977&\quad&2{,}4804090&
\end{array}\]}
\[\begin{array}{r@{\quad}r}
\text{Factor von }\theta\ldots&2{,}4804098\\[.4em]
\frac12\pi&0{,}1961199\\\cline{2-2}
\text{Erdquadrant}&2{,}6765297\\[.4em]
&7{,}0000000\\\cline{2-2}
\text{Factor Const.}&4{,}3234703\\[.4em]
&2{,}4811272\\\cline{2-2}
a&6{,}8045975
\end{array}\]
\noindent Also die Länge des Bogens
{\small
\[=\frac\theta{90^0}.10000000^{\mathrm m}+21060{,}580\cos T\sin T'+17{,}418\cos T'\sin T''.\]}
\NumberedItem{14}
\par\indent Die Winkel finden sich
\[\begin{array}{l|ll}
\sin2u=g\sin T&g\;\ldots\;8{,}9095925&h\;\ldots\;9{,}9992822\\[.4em]
&g'\;\ldots\;7{,}2185606&h'\;\ldots\;9{,}9999996\\[-.3em]
\sin T'=\dfrac{h\sin T}{\cos u^2}&g''\;\ldots\;3{,}8350620&h''\quad0
\end{array}\]
\clearpage\SourceLeaf{238}\ArithmeticHead{231}{R}{ZUR RECTIFICATION DER ELLIPSE.}
\noindent\begin{minipage}[t]{.33\linewidth}
\noindent Z.\ B.\ \(51.31.48=\varphi\)
\[\begin{array}{r}
0{,}09986\\[.4em]
00144\\\hline
\end{array}\]
\noindent \(38.33.46=T\)
\[\begin{array}{r@{\;}r}
\sin T&9{,}79475\\[.4em]
&8{,}90959\\[.4em]
&9{,}99928\\\cline{2-2}
&9{,}79403\\[.4em]
&28\\\cline{2-2}
\sin T'&9{,}79431\\[.4em]
T'=&38.31.0
\end{array}\]
\end{minipage}\hfill\vrule\hfill\begin{minipage}[t]{.62\linewidth}
\noindent Man führe noch ein
\[\begin{aligned}
a\operatorname{tg}U&=b\operatorname{tg}T\\[.4em]
a'\operatorname{tg}U'&=b'\operatorname{tg}T'\\[.4em]
a''\operatorname{tg}U''&=b''\operatorname{tg}T''.
\end{aligned}\]
\noindent Dann ist:
\[\sin T'=\frac{\sin\frac12(T+U)}{\cos\frac12(T-U)}\]
{\small
\[\begin{aligned}
\cos U'&=\frac{a\sin T'}{a'\operatorname{tg}T}=\frac{b\sin T'}{a'\operatorname{tg}U}\\[.6em]
&=\frac{b'\sin T'}{a'\surd\operatorname{tg}T.\operatorname{tg}U}\\[.6em]
\cos T'&=\frac{\sin U'}{\surd\operatorname{tg}T.\operatorname{tg}U}\\[.6em]
b'\cos T'&=\frac{a\sin U'}{\operatorname{tg}T}\\[.6em]
b'\cos T'&=\frac{b\sin U'}{\operatorname{tg}U}
\end{aligned}\]}
\end{minipage}
\[\begin{aligned}
\frac{\cos(T+U)}{\cos T\cos U}&=\frac{\cos(T'+U')\cos(T'-U')}{\cos{T'}^2}\\[.9em]
\frac{\cos(T+U)}{\sin T\sin U}&=\frac{\cos(T'+U')\cos(T'-U')}{\sin{U'}^2}\\[.9em]
\frac{2\cos(T+U)}{\surd\sin2T.\sin2U}&=\frac{\cos(T'+U').\cos(T'-U')}{\cos T'.\sin U'}\mathpunct{\raisebox{.5ex}{.}}
\end{aligned}\]
\clearpage\endgroup
